\newpage
\section*{Задание 6: Генерация технической документации}
\subsection{Постановка задачи}
\noindent
Алгоритм: Прямой проход нейронной сети (Forward Propagation) \\
Исходный код: оптимизированная реализация класса OptimizedNeuralNet (Python) из Задания 3.
Код не использует внешних библиотек (NumPy и т.п.), выполняет умножение матриц через списки и встроенные функции.\\
\newline
Задача: с использованием LLM (основная модель – gemini-3, контрольная – deepseek R1) сгенерировать техническую документацию для данного кода:
\begin{itemize}
	\item Встроенную документацию (docstrings) в стиле Google Style.
	\item Внешний файл README.md с описанием проекта, требованиями, примерами использования.
	\item Провести верификацию на галлюцинации и добавить теоретическое обоснование.
\end{itemize}
\subsection{Гипотеза и системный промпт}
\indent
Гипотеза:
Ролевая установка «опытный технический писатель, специализирующийся на документировании Python-кода для Data Science» в сочетании с явным указанием стандарта Google Style и требованием описывать типы, исключения и алгоритмическую сложность позволит получить качественную, свободную от галлюцинаций документацию.\\
\newline
\indent
Системный промпт:
\textit{
Ты – опытный технический писатель и Senior Python-разработчик. Твоя задача – создавать профессиональную документацию для кода на русском, предназначенного для задач машинного обучения.
Ты строго следуешь стандарту Google Style для docstrings. В каждом docstring ты обязательно указываешь:
\begin{itemize}
	\item Краткое описание (одна строка).
	\item Подробное описание (если необходимо).
	\item Args: имена, типы, смысл каждого параметра.
	\item Returns: тип и описание возвращаемого значения.
	\item Raises: какие исключения могут быть выброшены.
	\item Сложность (Big O) для времени и памяти (где уместно).
Ты не добавляешь вымышленных деталей и всегда сверяешься с предоставленным кодом.
Формат ответа – только код или текст документации, без лишних комментариев.
\end{itemize}
}
\subsection{Журнал диалога и итераций}
\noindent
\textbf{Итерация 1: Генерация docstrings (базовая)} \\
\indent
Промпт:
\textit{Вот код класса OptimizedNeuralNet (оптимизированная нейросеть прямого распространения). Сгенерируй для него docstrings в Google Style. Добавь описание для каждого метода (\_\_init\_\_, \_validate\_dimensions, forward). Не меняй реализацию.} \\
\newline
\indent
Ответ модели (Gemini-3): 
\lstinputlisting[language=Python]{z6/1-non-utf.py}
Анализ результата:
Сложность указана верно (умножение матриц через zip). Исключения не забыты. Теперь docstrings соответствуют промышленному стандарту. \\
\newline
\noindent
\textbf{Итерация 2: Генерация README.md} \\
\indent
Промпт:
\textit{
На основе кода и docstrings создай файл README.md в формате Markdown. Включи разделы:
\begin{itemize}
	\item Описание проекта: Требования (версия Python, зависимости – если нет внешних, так и укажи)
	\item Установка / запуск: Примеры использования (2-3 примера с конкретными входными данными и выводом)
	\item Структура проекта: Теоретическое обоснование (кратко, почему оптимизация через транспонирование)
\end{itemize}
}
\newpage
Ответ модели (Gemini-3): 
\lstinputlisting[language={}]{z6/2-non-utf.md}
Анализ результата:
README читаем, примеры синтаксически верны. Теоретическое обоснование етсь. Отсутствует раздел «Лицензия» (не критично).\\
Корректировка:
Попросить добавить предупреждение о возможных ошибках размерности. \\
\newline
\noindent
\textbf{Итерация 3: Верификация консистентности (самокоррекция)} \\
\indent
Промпт:
\textit{Выступи в роли критика. Сравни следующий код (класс  \\ OptimizedNeuralNet) и сгенерированную тобой документацию (docstrings + \\ + README). Найди несоответствия: описания параметров, типов, поведения. Если есть галлюцинации – укажи их.}
\newpage
\indent
Ответ модели (Gemini-3):
\begin{figure}[H]
	\centering
	\includegraphics[width=17cm]{z6/3-1.png}
	\caption{Итерация 3. ч.1}
\end{figure}
\newpage
\begin{figure}[H]
	\centering
	\includegraphics[width=17cm]{z6/3-2.png}
	\caption{Итерация 3. ч.2}
\end{figure}
\newpage
\begin{figure}[H]
	\centering
	\includegraphics[width=17cm]{z6/3-3.png}
	\caption{Итерация 3. ч.3}
\end{figure}
\noindent
Анализ результата:
Неверные типы аргументов и не все исключения указаны в функциях. Результат примера в README вымышлен. Нет защиты в неверных матриц в forward, а также лишние неиспользуемые переменные. \\
\newline
Корректировка:
Принять исправления, обновить README. \\
\newpage
\noindent
\textbf{Итерация 4: исправления} \\
\indent
\textit{
Код: Исправь типы в функциях чтобы они соответсвовали входным значениям и исправь Docstring к ним.
Добавь в forward проврку входного батча на размерность и пустоту.
Удали неиспользуемые переменные и импорты. Обнови для всего кода Docstring. \\
\newline
Документация (README): Вместо "придуманных"\ значений, расчитай результат модели прямого прохода в примере. \\
\newline
Ответ дай сначала блоком кода Python, а потом блоком кода Markdown (для README).
} \\
\newline
Ответ модели (Gemini-3, Код):
\lstinputlisting[language=Python]{z6/4-non-utf.py}
\newpage
Ответ модели (Gemini-3, Документация):
\lstinputlisting[language={}]{z6/4-non-utf.md}
\newpage
\noindent
\textbf{Итерация 5: Добавление полного теоретического раздела} \\
\indent
Промпт:
\textit{Добавь в README отдельный раздел «Теоретическое обоснование» объёмом 3–5 абзацев. Опиши: почему выбран рекуррентный обход через zip вместо вложенных индексов, как транспонирование снижает cache-miss, и как оценивалась сложность.} \\
\newline
Ответ модели (Gemini-3):
\lstinputlisting[language={}]{z6/5-non-utf.md}
\subsection{Артефакты}
Итоговый код с docstrings: (Итерация 4, код Python) \\
\indent
Файл README\.md: (Итерация 5, код Markdown)
\subsection{Сравнение моделей}
Финальный промпт (объединяющий все этапы) был подан контрольной модели DeepSeek R1 без изменений.
Результаты сравнения:
\begin{table}[h!]
	\centering
	\begin{tabular}{|l|p{5cm}|p{5cm}|}
		\hline
		\textbf{Критерий} & \textbf{модель (a)} & \textbf{модель (b)} \\
		\hline
		Соответствие стандартам  \\ (docstrings / README) & 5 & 4 \\
		\hline
		Полнота документации  & 5 & 4\\
		\hline
		Отсутствие галлюцинаций & 4 & 4 \\
		\hline
		Качество объяснений & 4 & 5 \\
		\hline
		\textbf{Итог} & 4.5 & 4.25\\
		\hline
	\end{tabular}
\end{table} \\
Итог: gemini-3 справилась лучше благодаря способности удерживать сложный контекст и следовать многошаговым инструкциям. deepseek R1 генерирует более краткие ответы, что в данном задании привело к потере важных разделов документации.
\newpage
\subsection{Вывод}
В ходе выполнения задания наиболее эффективными техниками промпт-инжиниринга оказались:
\begin{itemize}
	\item Ролевое моделирование (явное указание «технический писатель, следующий Google Style») – обеспечило корректный формат docstrings.
	\item Итеративное уточнение – позволило добавить пропущенные разделы (сложность, исключения).
	\item Самокоррекция (верификация) – выявила мелкие недочёты в README, которые были исправлены без ручного редактирования.
	\item Обогащение контекста (передача полного кода и предыдущих ответов) – предотвратило галлюцинации.
\end{itemize}
Основная модель (gemini-3) продемонстрировала высокую точность следования инструкциям и способность генерировать структурированную техническую документацию, близкую к человеческой. Контрольная модель (deepseek R1) дала рабочий, но менее детальный результат, что подтверждает необходимость тщательной настройки промптов под конкретную LLM.